\documentclass[11pt]{article}

\usepackage[margin=1.2in]{geometry}
\usepackage{amsmath}
\usepackage{amssymb}
\usepackage{amsthm}

\newtheorem{theorem}{Theorem}[section]
\newtheorem{lemma}[theorem]{Lemma}
\theoremstyle{definition}
\newtheorem{example}[theorem]{Example}
\theoremstyle{remark}
\newtheorem{remark}[theorem]{Remark}

\newcommand{\id}{\mathrm{id}}
\newcommand{\Fix}{\mathrm{Fix}}

\title{A Self-Referential Valuation No-Go and the\\ Forced Symmetry-Breaking of the Residual}
\author{Joe Hernandez\thanks{%% JOE: fill in email and affiliation (or independent-researcher line) before posting.
Email and affiliation to be added.}}
\date{July 2026}

\begin{document}

\maketitle

\begin{abstract}
We prove a short structural theorem about self-referential observers carrying a two-valued admissibility grading. The observer is modelled as a point of an object $A$ in a category with finite products; valuations are morphisms $p \colon A \to B$ into a two-element grading object; the grading symmetry is a fixpoint-free involution $\alpha$ on $B$. Two facts follow. (I)~No weakly point-surjective $T \colon A \times A \to B$ exists (a Cantor--Lawvere diagonal obstruction), and, separately, no total valuation is $\alpha$-invariant: any valuation the observer commits to breaks the grading symmetry. (II)~Under a purely group-theoretic, non-circular arena/value partition, the residual is provably a value. We locate the result relative to Lawvere/Cantor fixed-point arguments and Curie-type invariant dichotomies; the contribution is the synthesis and the non-circular partition, not a new fixed-point theorem. The main lemma and both corollaries are machine-checked in Lean~4.
\end{abstract}

\section{Informal statement}
\label{sec:informal}

This is a short, self-contained structural theorem. It requires no physics and no prior context: the reader needs only elementary category theory (finite products, the diagonal) and elementary group theory (a group action and its fixed points). The paper separates six parts explicitly -- the setting (Section~\ref{sec:setting}), the assumptions (Section~\ref{sec:assumptions}), the theorem (Section~\ref{sec:theorem}), the proof (Section~\ref{sec:proof}), examples and counterexamples (Section~\ref{sec:examples}), and the limits of the result (Section~\ref{sec:limits}) -- and adds a dedicated novelty map (Section~\ref{sec:novelty}), which is where the honest weight of the contribution is assessed.

A self-referential \emph{observer} is modelled as a point of a structure that also carries a two-valued \emph{admissibility grading} on that same structure. The grading comes with an involution $\alpha$ that exchanges the two grades (``admissible'' and ``inadmissible''). We assume this involution is \emph{fixpoint-free} (nothing is both admissible and inadmissible). Two elementary facts follow:

\begin{itemize}
\item \textbf{(I) No self-referential closure, and no symmetric total valuation.} There is no self-referential total representation of all admissibility valuations (a Cantor/Lawvere diagonal obstruction); and, separately, no total admissibility valuation is $\alpha$-invariant. Hence any total valuation the observer actually commits to must \emph{break} the grading symmetry.
\item \textbf{(II) The residual is necessarily a ``value,'' in a non-circular sense.} Define, purely group-theoretically: \emph{arena} $=$ a quantity invariant under the acting symmetry; \emph{value} $=$ a quantity whose specification requires breaking that symmetry. This partition is decidable from the group action alone, with no reference to what is ``forced'' or ``selected'' by any process. Under this definition the residual forced by (I) is provably a value.
\end{itemize}

The synthesis packages these as: \emph{the selection an observer is forced to make is, provably and non-circularly, a symmetry-breaking one.} One caution up front, so the reader is not misled about the mathematical weight: the symmetry-breaking conclusion rests entirely on the \emph{elementary} fixed-point-count lemma (I-b) / Lemma~\ref{lem:C} -- ``a nonempty domain cannot map into an empty fixed-set'' -- and \emph{not} on Lawvere's theorem. Lawvere / (I-a) contributes the \emph{independent} ``no self-enumeration'' fact; it does not underwrite the symmetry-breaking result. We are careful below to state exactly which pieces are classical and which is the new packaging.

\begin{remark}[On ``forced'' and ``commit'': informal language]
\label{rem:forced}
Throughout, the words ``forced,'' ``commit,'' and ``the observer commits to a valuation'' are an \emph{informal} motivating gloss, not a formal predicate. No object in the formal development is the referent of ``forced'': the residual $d = \alpha \circ T \circ \Delta$ (notation of Section~\ref{sec:setting}) depends on a \emph{choice} of $T$, and there is no theorem that the observer must adopt any particular valuation. What the theorem actually proves is the non-existence of an invariant or self-enumerating valuation (parts I-a and I-b), plus the group-theoretic character of any total valuation (part II). One could formalize ``commitment'' trivially as ``existence of some total $p \colon A \to B$'' (immediate for nonempty $A$), but nothing below depends on doing so; read ``forced'' as a reading, not a proved modality.
\end{remark}

\section{The abstract setting (Part 1)}
\label{sec:setting}

Work in a category $\mathcal{C}$ with finite products (the reader may take $\mathcal{C} = \mathbf{Set}$; nothing below uses more). Write $1$ for the terminal object and call a morphism $x \colon 1 \to X$ a \emph{point} of $X$. In $\mathbf{Set}$, points of $X$ are its elements, and we use element notation freely.

\paragraph{Objects.}
\begin{itemize}
\item $A$ -- the \emph{arena of observer-states}. Its points are the states the observer can occupy. The observer is self-referential in the sense made precise below: a state of $A$ can serve both as a code (an admissibility rule) and as the argument that rule is applied to.
\item $B$ -- the \emph{value object}, a two-element object $B = \{0, 1\}$. Read $0 =$ admissible, $1 =$ inadmissible. $B$ is the object in which gradings and valuations take their values.
\end{itemize}

\paragraph{The diagonal.} Finite products give the diagonal $\Delta \colon A \to A \times A$, $\Delta(a) = (a, a)$. This is the entire formal content of ``self-reference'': a state applied to itself.

\paragraph{Valuations.} A \emph{valuation} is a morphism $p \colon A \to B$; it grades each state as admissible or inadmissible.

\paragraph{The grading involution (the firewall).} A morphism $\alpha \colon B \to B$ with $\alpha \circ \alpha = \id_B$. The intended $\alpha$ is the \emph{swap} $\alpha(0) = 1$, $\alpha(1) = 0$ -- the map that exchanges admissible and inadmissible. $\alpha$ acts on valuations by post-composition: $(\alpha \circ p)(a) = \alpha(p(a))$.

\paragraph{Self-referential closure.} A morphism $T \colon A \times A \to B$ is a candidate \emph{closure}: its ``rows'' $T(a_0, -) \colon A \to B$ are meant to enumerate valuations, with $a_0$ the code and the second argument the input. $T$ is \emph{weakly point-surjective} if every valuation is some row: for each $p \colon A \to B$ there is a point $a_0 \colon 1 \to A$ such that $T(a_0, x) = p(x)$ for all points $x \colon 1 \to A$. A weakly point-surjective $T$ is exactly a self-referential structure that represents \emph{all} of its own valuations internally.

That is the whole setting: a self-referential arena $A$, a two-valued grading object $B$ with an involution $\alpha$, and the notion of a total internal representation $T$ of $A$'s valuations.

\section{The assumptions (Part 2)}
\label{sec:assumptions}

\begin{itemize}
\item \textbf{(A1) Finite products.} $\mathcal{C}$ has finite products (hence the diagonal $\Delta$ and the pairing used in the proof). Satisfied by any cartesian category, in particular $\mathbf{Set}$.
\item \textbf{(A2) Two-valued grading (reading convention).} $B$ has \emph{exactly} the two distinct points $0, 1$ ($|B| = 2$) that $\alpha$ interchanges; the grading is genuinely two-valued, not collapsed to one point. This is a reading convention that fixes the intended model; the load-bearing content is carried entirely by A3. (If one prefers to allow $|B| > 2$, replace A2 by the requirement that $\alpha$ be fixpoint-free on \emph{all} of $B$, i.e.\ strengthen A3 to cover the extra points; the results are unchanged, and a third \emph{fixed} grade dissolves the theorem -- Example~\ref{ex:thirdgrade}.)
\item \textbf{(A3) Fixpoint-freeness of the firewall.} $\alpha$ has \emph{no fixed point}: there is no point $b \colon 1 \to B$ with $\alpha \circ b = b$. For the swap on $\{0, 1\}$ this is immediate (admissible $\neq$ inadmissible). This is the \emph{sole load-bearing hypothesis}: every conclusion (I-a) and (I-b) is a strict consequence of A3 (given nonempty $A$ in A4), and $|B| = 2$ is inessential to the proofs.
\item \textbf{(A4) Nonempty arena} (used only for the invariant-valuation half): $A$ has at least one point.
\end{itemize}

No other assumption is used. In particular no topology, measure, dynamics, order, or physical structure enters.

\section{The theorem (Part 3)}
\label{sec:theorem}

\begin{theorem}
\label{thm:main}
Let $\mathcal{C}$ be a category with finite products, $A$ a nonempty object, $B$ a two-valued object, and $\alpha \colon B \to B$ a fixpoint-free involution (A1--A4). Then:

\emph{(I-a) No closure.} No weakly point-surjective $T \colon A \times A \to B$ exists. Concretely, for every $T \colon A \times A \to B$ the \emph{diagonal valuation} $d = \alpha \circ T \circ \Delta \colon A \to B$ is not a row of $T$; it is a valuation that $A$ cannot represent from inside.

\emph{(I-b) No invariant valuation.} No valuation $p \colon A \to B$ is $\alpha$-invariant: $\alpha \circ p = p$ is impossible for nonempty $A$. Equivalently, every total valuation the observer commits to satisfies $\alpha \circ p \neq p$ -- it breaks the grading symmetry.

\emph{(II) The residual is a value.} Let a group $G$ act on $B$, and extend the action to valuations by post-composition. Call a valuation \emph{arena-type} if it is $G$-invariant and \emph{value-type} if it is not. This dichotomy is decided by the $G$-action alone. Taking $G = \langle\alpha\rangle \cong \mathbb{Z}/2$, part (I-b) says every total valuation is value-type. Hence the selection an observer is forced into by (I) is necessarily a \emph{value} (a symmetry-breaking selection), and this is a group-theoretic fact, not a definitional restatement of ``forced.''
\end{theorem}

Parts (I-a) and (I-b) are two distinct elementary consequences of A3 (they are not literally the same statement: (I-a) is about no total \emph{enumeration}, (I-b) about no single \emph{invariant} valuation). Part (II) supplies the non-circular reading under which the forced residual is a genuine value.

\section{The proof (Part 4)}
\label{sec:proof}

The proof is purely mathematical and uses no numerical input. (Finite instances have been machine-checked as confirmation only; see Remark~\ref{rem:machine} in Section~\ref{sec:examples}. They are not part of the proof.)

\subsection{Lawvere's weak fixed-point lemma}

\begin{lemma}[Lawvere \cite{lawvere1969}; weak form, Yanofsky \cite{yanofsky2003}]
\label{lem:L}
If there is a weakly point-surjective $T \colon A \times A \to B$, then every endomorphism $\alpha \colon B \to B$ has a fixed point: a point $b \colon 1 \to B$ with $\alpha \circ b = b$.
\end{lemma}

\begin{proof}
Given $\alpha$, form the valuation $f = \alpha \circ T \circ \Delta \colon A \to B$, i.e.\ $f(a) = \alpha(T(a, a))$. By weak point-surjectivity there is a point $a_0$ with $T(a_0, x) = f(x)$ for all points $x$. Evaluate at $x = a_0$:
\[
T(a_0, a_0) = f(a_0) = \alpha(T(a_0, a_0)).
\]
So $b := T(a_0, a_0)$ satisfies $\alpha \circ b = b$.
\end{proof}

\begin{proof}[Proof of (I-a)]
Contrapositive of Lemma~\ref{lem:L}. By A3, $\alpha$ has no fixed point; therefore no weakly point-surjective $T$ exists. Moreover the specific valuation exhibited in the proof, $d = \alpha \circ T \circ \Delta$, cannot be a row of $T$: if $d = T(a_0, -)$ for some code $a_0$, then evaluating at $a_0$ gives $T(a_0, a_0) = d(a_0) = \alpha(T(a_0, a_0))$, a fixed point of $\alpha$, contradicting A3. So $d$ is the concrete unrepresented (residual) valuation.
\end{proof}

\subsection{The fixed-point-count corollary}

\begin{lemma}[Elementary fixed-point count]
\label{lem:C}
Let $\alpha \colon B \to B$ and $p \colon A \to B$. If $\alpha \circ p = p$ then $p$ factors through the fixed-point subobject $\Fix(\alpha) = \{b : \alpha(b) = b\}$ (the equalizer of $\alpha$ and $\id_B$). If $\Fix(\alpha)$ is empty and $A$ is nonempty, then $\alpha \circ p = p$ is impossible.
\end{lemma}

\begin{proof}
$\alpha \circ p = p$ says every value $p(a)$ satisfies $\alpha(p(a)) = p(a)$, i.e.\ $p(a) \in \Fix(\alpha)$; so $p$ factors through $\Fix(\alpha)$. If $\Fix(\alpha) = \emptyset$ and $A$ has a point $a$, then $p(a)$ would be a point of the empty object, which is impossible.
\end{proof}

\begin{proof}[Proof of (I-b)]
By A3, $\Fix(\alpha) = \emptyset$; by A4, $A$ is nonempty. Lemma~\ref{lem:C} gives $\alpha \circ p \neq p$ for every valuation $p$.
\end{proof}

\subsection{The non-circular partition and the identification (II)}
\label{sec:partition}

Let a group $G$ act on $B$; extend to valuations $p \colon A \to B$ by $(g \cdot p)(a) = g \cdot (p(a))$. Define a valuation to be \emph{arena-type} if $g \cdot p = p$ for all $g \in G$ (invariant) and \emph{value-type} otherwise.

\paragraph{Non-circularity.} Whether a given $p$ is arena-type or value-type is determined entirely by the $G$-action on $B$ and the values of $p$ -- a computation in group theory. It makes no reference to any external notion of a quantity being ``forced,'' ``selected,'' ``dynamically determined,'' or ``chosen.'' In particular the statement ``the arena-type quantities are exactly the forced ones'' is \emph{synthetic}: it relates two independently-defined predicates (invariance on one side, forcing on the other) and could a priori be false. This is what saves the partition from the tautology ``arena $:=$ the forced things, value $:=$ the unforced things.'' (See Example~\ref{ex:partition} for the logical independence made concrete.)

\paragraph{The identification.}
Take $G = \langle\alpha\rangle \cong \mathbb{Z}/2$, generated by the fixpoint-free involution. A valuation is arena-type iff $\alpha \circ p = p$. By (I-b) no total valuation is arena-type; every total valuation is value-type. Therefore the residual valuation the observer is forced to commit to (its existence and unrepresentability supplied by (I-a); its non-invariance by (I-b)) is a \emph{value} in the group-theoretic sense. \qed

\begin{remark}[One-directional vacuity: honest disclosure]
\label{rem:vacuity}
The implication ``forced $\Rightarrow$ value'' is \emph{vacuous} in that direction: by (I-b) \emph{every} total valuation is value-type, so learning that the residual is a value says nothing specific about the residual -- it is a property shared by all total valuations. The substantive content of Part~II is therefore \emph{not} ``forced $\Rightarrow$ value'' but the two facts that survive this vacuity: (i) the purely negative no-go \emph{no $\alpha$-invariant total valuation exists} (I-b), and (ii) the non-circularity of the arena/value partition (the ``invariant iff forced'' reading is synthetic, not definitional; see the non-circularity paragraph above and Example~\ref{ex:partition}). The information-carrying direction is the \emph{absence} of any invariant residual, not the presence of a value.
\end{remark}

\section{Examples and counterexamples (Part 5)}
\label{sec:examples}

\begin{example}[The generating instance is Cantor]
\label{ex:cantor}
Take $A$ any set, $B = \{0, 1\}$, $\alpha =$ swap. A weakly point-surjective $T \colon A \times A \to B$ is exactly an enumeration of $\mathcal{P}(A)$ by $A$ (a valuation $A \to B$ is the indicator of a subset). Part (I-a) then reads: no such enumeration exists, and, writing $S_a \subseteq A$ for the subset whose indicator is the row $T(a, -)$, the unrepresented $d = \alpha \circ T \circ \Delta$ is the indicator of the Cantor diagonal set $\{a : a \notin S_a\}$ \cite{cantor1891}. So with two grades and the swap, Part~I \emph{is} Cantor's theorem. This is stated up front because it fixes what is and is not new (Section~\ref{sec:novelty}).
\end{example}

\begin{example}[Control: a firewall with a fixed point dissolves the theorem]
\label{ex:control}
Drop A3: let $\alpha = \id_B$ (or any $\alpha$ with a fixed grade). Then Lemma~\ref{lem:L} has a trivial fixed point, weakly point-surjective $T$ can exist, and constant valuations are $\alpha$-invariant (arena-type). No selection is forced. This shows the conclusion is contributed by fixpoint-freeness specifically, not by ``a residual exists'' in general. The hypothesis A3 is load-bearing, not decorative.
\end{example}

\begin{example}[Counterexample: a genuine third grade breaks the argument]
\label{ex:thirdgrade}
Let $B = \{\text{below}, \text{boundary}, \text{above}\}$ and let $\alpha$ exchange below and above but fix boundary. Now $\alpha$ has a fixed point, Lemma~\ref{lem:L} applies with $b = \text{boundary}$, and an $\alpha$-invariant valuation (constantly boundary) exists. So a genuine neutral middle grade defeats both (I-a) and (I-b). The theorem needs the grading to be \emph{effectively two-valued with no fixed grade} (A2 $+$ A3); any admissibility scheme with a stable ``neutral'' verdict is outside its scope. This is the sharpest limit on the setting.
\end{example}

\begin{example}[The partition is substantive, not circular]
\label{ex:partition}
Consider $G = \mathbb{Z}/2$ acting on $B = \{0, 1\}$ by swap. The constant-pair feature ``the unordered set $\{0, 1\}$'' is $G$-invariant (arena-type); a specific choice of element $0$ or $1$ is not $G$-invariant (value-type). Logical independence from ``forcing'' is visible: one can specify an action under which a quantity is invariant while a \emph{separate} dynamics fails to determine it (an unforced invariant), or is non-invariant while a dynamics does determine it (a forced value). Because both mismatches are describable, ``invariant iff forced'' is a claim with content that could fail -- exactly the mark of a non-vacuous partition. (In the source research program this was tested against a battery of concrete quantities; here we only need the logical point.)
\end{example}

\begin{example}[Positive existence when hypotheses are met]
\label{ex:positive}
For $A$ nonempty and $\alpha$ the swap, (I-a) and (I-b) both hold, and the concrete unrepresented residual $d = \alpha \circ T \circ \Delta$ exists for every candidate $T$. So the theorem is not vacuously true by absence of models: the setting is inhabited and the residual is exhibited, not merely asserted.
\end{example}

\begin{remark}[Machine-checked formalization: confirmation, not proof]
\label{rem:machine}
The main lemma (Lemma~\ref{lem:L} together with its contrapositive engine) and both corollaries -- the no-closure statement (I-a), including the concrete ``residual is not a row'' form, and the no-invariant-valuation statement (I-b) -- have been formalized in Lean~4 and kernel-checked, with no unproven placeholders; finite-instance numerical confirmations (exhaustive over all $T$ for $|A| \in \{1, 2, 3\}$, together with the Cantor cross-check, the fixpoint-freeness check, and the control cases of Examples~\ref{ex:control} and~\ref{ex:thirdgrade}) accompany the formalization. Both are available in the public repository at \texttt{https://github.com/disruptionjoe/gu-formalization}. These runs confirm the elementary lemmas on small cases; they are evidence, not premises. The proof stands on Lemmas~\ref{lem:L} and~\ref{lem:C} alone.
\end{remark}

\section{Novelty map (dedicated assessment)}
\label{sec:novelty}

This section maps the result against the nearest prior art and grades its novelty honestly. The three possible verdicts are: \textbf{(a) genuinely novel} (a real synthesis with a new load-bearing statement); \textbf{(b) novel packaging} (correct and clearly stated, but the pieces are known, the value being in the framing); \textbf{(c) subsumed} (already in the literature under another name).

\subsection{Lawvere's fixed-point theorem; Yanofsky}

Part (I-a) \emph{is} Lawvere's weak fixed-point theorem \cite{lawvere1969}, in the elementary form popularized by Yanofsky \cite{yanofsky2003}, applied to a two-element $B$ with a fixpoint-free $\alpha$. Yanofsky's ``universal approach'' already unifies Cantor, Russell, G\"odel, Tarski, and the halting problem as instances of exactly this scheme. The fixed-point core is therefore \emph{not new}; with two grades and the swap it is literally Cantor \cite{cantor1891}. What is not in Lawvere/Yanofsky is the \emph{framing}: the two-element $B$ read as an admissibility firewall carried by a self-referential observer, and the identification (via Part~II) of the forced residual as a symmetry-breaking selection. That framing is additive, not a strengthening of the theorem.

\subsection{Kochen--Specker and contextuality}

The Kochen--Specker theorem \cite{kochenspecker1967} is also a two-valued valuation no-go: for a Hilbert space of dimension $\geq 3$ there is no global $\{0, 1\}$-valuation on the projection lattice compatible with the functional (FUNC) constraints. Our result is \emph{also} a two-valued valuation no-go with a forced residual, so the genus is shared. But the \emph{mechanism} differs and they are not the same theorem: Kochen--Specker's obstruction is geometric/combinatorial (a Gleason-type coloring impossibility forced by the orthogonality structure of quantum observables in dimension $\geq 3$), and its ``context'' is a maximal commuting set. Ours is a \emph{self-referential diagonal} obstruction forced by fixpoint-freeness of a grading involution, with self-application (the diagonal $\Delta$) as the essential ingredient; dimension, orthogonality, and measurement contexts play no role. Abramsky and Brandenburger \cite{abramskybrandenburger2011} recast contextuality sheaf-theoretically (no global section of a presheaf of local valuations); that is closer in spirit (a global-vs-local obstruction) but again the obstruction is measurement-context gluing, not self-reference. Conway--Kochen \cite{conwaykochen2006} (free will) and Bell \cite{bell1966} add physical (locality/determinism) premises we do not use. Verdict on this axis: \emph{related but distinct}; ours is not Kochen--Specker in disguise, though both belong to the family ``no consistent global two-valuation.''

\subsection{Curie's principle}

Part~II (``value $=$ requires symmetry-breaking to specify'') is the elementary invariant / symmetry-breaking-order-parameter dichotomy -- a quantity is arena-type iff it is $G$-invariant and value-type (an order parameter) otherwise. This is Curie-\emph{flavored}, but it is \emph{not} Curie's principle proper: Curie's principle \cite{curie1894} is a \emph{causal} claim (symmetric causes have symmetric effects, so an asymmetry of effects requires an asymmetry of causes), whereas the dichotomy here is the textbook invariant-vs-order-parameter distinction, with no cause and no effect in the definition. Earman \cite{earman2004} and Norton's ``Curie's Truism'' \cite{norton2016} both note that precise formulations of Curie's principle tend toward the \emph{analytic} (nearly vacuous in quantum field theory). We do \emph{not} claim the \emph{forcing} rescues Curie from that analyticity worry: as (I-b) shows, every total valuation is value-type, so ``the forced residual is a value'' is vacuous in that direction (see Remark~\ref{rem:vacuity} and Section~\ref{sec:limits}). The non-vacuous content is the purely negative fact -- \emph{no $\alpha$-invariant total valuation exists} -- together with the non-circular group-theoretic partition (Section~\ref{sec:partition}, Example~\ref{ex:partition}). That negative fact is exactly the analytic-leaning fact Earman warns about; our contribution is to tie it to a \emph{valuation no-go} (Part~I) and to a non-circular partition, not to escape analyticity via ``forcing.''

\subsection{The general Lawvere-paradox literature}

Yanofsky \cite{yanofsky2003} and the broader categorical-diagonal tradition already subsume the fixed-point machinery and many of its instances. The present result is not a new theorem \emph{in that tradition}; it is an application of it to a graded/firewall setting joined to a Curie-style partition.

\subsection{Honest verdict: (b) novel packaging}

The load-bearing mathematical facts (Lemmas~\ref{lem:L} and~\ref{lem:C}) are classical; with two grades Part~I is Cantor and in general it is Lawvere/Yanofsky, and Part~II is the elementary Curie-flavored invariant / order-parameter dichotomy (not Curie's causal principle). What is genuinely contributed, and is not simply any one of those pieces, is the \emph{synthesis}: (i) reading the two-element grading as an admissibility firewall on a self-referential observer, so that the classical diagonal fires on ``the observer's own admissibility valuation''; (ii) proving the forced residual is not merely unforced but \emph{symmetry-breaking}, via the fixpoint-count corollary; and (iii) supplying a non-circular, group-theoretic arena/value partition that makes ``the residual is a value'' a substantive rather than definitional claim, and that ties the Curie-flavored invariant / order-parameter dichotomy to a valuation no-go. This is a correct, clearly-stated, well-mapped synthesis whose parts are known -- the respectable and common (b) outcome. The closest single prior result is Yanofsky \cite{yanofsky2003} for Part~I and Earman's analysis of Curie's principle \cite{earman2004} for Part~II. We do \emph{not} claim (a): no individual step is a new theorem, and inflating the synthesis to ``genuinely novel'' would not survive a specialist who recognizes the Cantor/Lawvere core. We also do not accept (c): no prior source states the combined ``forced residual is provably a symmetry-breaking value, under a non-circular partition'' as a single result, so it is not strictly subsumed.

\section{What the theorem does NOT establish (Part 6)}
\label{sec:limits}

\begin{itemize}
\item \textbf{No physical or operator-algebra realization.} The theorem is purely structural. A separate attempt to realize $A$ as the observable algebra of a spacetime region and $\alpha$ as a grading of a modular conjugation -- so that the abstract firewall would become a physical one -- \emph{did not go through}: in the interacting continuum the required per-region construction fails to exist and to cohere across overlaps (this break is recorded in the accompanying repository). Nothing in this paper depends on that realization, and none of it is claimed here. The abstract theorem stands; the physical instantiation does not.
\item \textbf{No dependence on, or support for, any physical theory.} The setting uses no physics. The theorem neither assumes nor provides evidence for any specific unified theory, field content, or model. Any physical theory can at most be \emph{one instance} of the abstract setting, and being an instance confers no confirmation on the theory.
\item \textbf{No specific value, constant, or count.} The theorem says the forced residual \emph{is a symmetry-breaking selection}; it does not compute which selection, nor any number, ratio, dimension, or count. It is an existence-and-character result, not a quantitative one.
\item \textbf{Not a physical prediction.} It is a logical/structural theorem about self-referential valuations. It makes no empirically testable claim on its own.
\item \textbf{Genericity, not uniqueness.} Because Part~II is the elementary invariant / order-parameter dichotomy (Curie-flavored, not Curie's causal principle), the arena/value partition it uses is \emph{generic} to symmetry-based structures, not a signature of any particular system. Genericity is not vacuity (the partition is still non-circular and falsifiable), but it means the theorem cannot be used to argue that any specific framework is distinguished.
\item \textbf{Novelty caveats (from Section~\ref{sec:novelty}).} Part~I is classical (Cantor / Lawvere / Yanofsky) and Part~II is the elementary Curie-flavored invariant / order-parameter dichotomy (not Curie's causal principle). The honest novelty grade is (b) novel packaging; the value is in correctness, self-containment, and the mapped synthesis, not in a new fixed-point core.
\item \textbf{Scope of the grading.} The theorem requires an effectively two-valued grading with a fixpoint-free involution (A2, A3). It says nothing about admissibility schemes with a stable neutral verdict; a genuine third fixed grade dissolves it (Example~\ref{ex:thirdgrade}).
\end{itemize}

\begin{thebibliography}{99}

\bibitem{lawvere1969}
F.~W. Lawvere,
\emph{Diagonal arguments and cartesian closed categories},
in: Category Theory, Homology Theory and their Applications II,
Lecture Notes in Mathematics, vol.~92, Springer, 1969, pp.~134--145.
Reprinted with author commentary in: Reprints in Theory and Applications of Categories, no.~15 (2006), pp.~1--13.

\bibitem{yanofsky2003}
N.~S. Yanofsky,
\emph{A universal approach to self-referential paradoxes, incompleteness and fixed points},
Bulletin of Symbolic Logic \textbf{9} (2003), no.~3, 362--386.

\bibitem{cantor1891}
G.~Cantor,
\emph{\"Uber eine elementare Frage der Mannigfaltigkeitslehre},
Jahresbericht der Deutschen Mathematiker-Vereinigung \textbf{1} (1891), 75--78.

\bibitem{kochenspecker1967}
S.~Kochen and E.~P. Specker,
\emph{The problem of hidden variables in quantum mechanics},
Journal of Mathematics and Mechanics \textbf{17} (1967), 59--87.

\bibitem{bell1966}
J.~S. Bell,
\emph{On the problem of hidden variables in quantum mechanics},
Reviews of Modern Physics \textbf{38} (1966), 447--452.

\bibitem{abramskybrandenburger2011}
S.~Abramsky and A.~Brandenburger,
\emph{The sheaf-theoretic structure of non-locality and contextuality},
New Journal of Physics \textbf{13} (2011), 113036.

\bibitem{conwaykochen2006}
J.~Conway and S.~Kochen,
\emph{The free will theorem},
Foundations of Physics \textbf{36} (2006), 1441--1473.

\bibitem{curie1894}
P.~Curie,
\emph{Sur la sym\'etrie dans les ph\'enom\`enes physiques, sym\'etrie d'un champ \'electrique et d'un champ magn\'etique},
Journal de Physique Th\'eorique et Appliqu\'ee \textbf{3} (1894), 393--415.

\bibitem{earman2004}
J.~Earman,
\emph{Curie's Principle and spontaneous symmetry breaking},
International Studies in the Philosophy of Science \textbf{18} (2004), no.~2--3, 173--198.

\bibitem{norton2016}
J.~D. Norton,
\emph{Curie's Truism},
Philosophy of Science \textbf{83} (2016), no.~5, 1014--1026.

\end{thebibliography}

\end{document}
